\scanpage{46}
Notons que pour deux espaces vectoriels $E,F$, on a un isomorphisme canonique
\[
(2)\qquad\overline{E\otimes F}=\overline E\otimes\overline F.
\]
En effet, l'application $(\overline x,\overline y)\longrightarrow\unsure{\overline{x\otimes y}|\overline{\overline x\otimes\overline y}}$ de $\overline E\unsure{\times|\otimes}\overline F$ dans $\overline{E\otimes F}$ est manifestement bilinéaire, donc se prolonge en une application linéaire $\overline E\otimes\overline F\longrightarrow\overline{E\otimes F}$
\margin{\deleted{la forme linéaire \par \unsure{application} \deleted{\unsure{linéaire}} antilinéaire. \par \unsure{norme}. Par suite, on a}}

Notons maintenant que si $E$ est un espace préhilbertien, alors la fonction $\overline x,\overline y\longrightarrow\unsure{(y,x)}$ sur $\overline E\times\overline E$ est une forme sesquilinéaire ; qu'on notera $(\overline x,\overline y)$ ; qui \deleted{\unsure{fait} \illegible{2} $\overline E$ \illegible{3}} est \deleted{\unsure{définie}} \inserted{\unsure{encore}} positive, car $(\overline x,\overline x)=(x,x)$.
Ceci permet de considérer $\overline E$ comme un espace préhilbertien.
\[
(3)\qquad(\overline x,\overline y)=(y,x)
\]
Comme on a $\|x\|=\|\overline x\|$, la norme de $\overline E$ coïncide avec celle de $E$. Si donc $E$ est séparé, $\overline E$ l'est, et \uline{si $E$ est un espace de Hilbert, $\overline E$ est aussi un espace de Hilbert}.
Dans ce dernier cas, on sait que l'application qui à tout $y\in E$ fait correspondre la forme $x\longrightarrow(x,y)$ sur $E$, est une \deleted{\unsure{isométrie}} \inserted{\unsure{appl. bijective}} \uline{antilinéaire} de $E$ \uline{sur} le dual $E'$, conservant la norme, par suite, on a \deleted{\unsure{de façon}} \inserted{ainsi} un isomorphisme canonique d'espaces normés :
\[
(4)\qquad E'=\overline E\qquad(E,\text{ espace de Hilbert})
\]
Cet isomorphisme permet de considérer $E'$ comme \deleted{\unsure{muni de la} \illegible{2}} \deleted{\unsure{forme de} \illegible{2}} un espace de Hilbert.
\mspage{47}
